% Résultats RPEM en simulation (V_c et f_z), découpés en trois parties (frame avec \only<1-3>)
%   partie 1 : V_c, f_z + coefficients a_w ... f_w
%   partie 2 : V_c, f_z + K_f, K_c, K_m et les efforts associés
%   partie 3 : V_c, f_z + W_c,max et critère J
% Les sous-figures sont découpées dans images/rpem_simu_fz_Vc.pdf_tex (grille 4 lignes x 4 colonnes).
% Vert = validé, orange = biaisé mais exploitable (a_w, b_w, c_w, d_w, f_w),
% rouge = non exploitable (K_f, K_m, F_er, M_P,ez).
% Les étiquettes verticales de K_f, K_c, K_m et des efforts sont réécrites (elles se chevauchaient dans l'export MATLAB).
% La macro \rpempanneau est définie dans le préambule de presentation.tex.
\providecommand{\rpemlargeur}{0.83\linewidth}
\sbox{\rpemboite}{\fontsize{3.5pt}{5pt}\selectfont\def\svgwidth{\rpemlargeur}\import{images/}{rpem_simu_fz_Vc.pdf_tex}}%
\begin{tikzpicture}[x=\rpemlargeur, y=\rpemlargeur]
    % bornes des cellules dans la figure d'origine (fraction de la largeur)
    \def\cA{0}\def\cB{0.252}\def\cC{0.4975}\def\cD{0.755}\def\cE{1}
    \def\rA{0}\def\rB{0.161}\def\rC{0.3132}\def\rD{0.4654}\def\rE{0.6175}
    % réglages du découpage propres à cette figure (lignes plus serrées)
    \def\rpemcoin{0.03}\def\rpemdeb{0.008}\def\rpemexA{0.035}\def\rpemexB{0.105}\def\rpemexH{0.007}\def\rpemexHb{0.0079}
    % disposition : V_c et f_z agrandis à gauche, panneaux de chaque partie à droite
    % (tous les panneaux à la même échelle : les deux côtés ont la même hauteur)
    \def\pas{0.258}\def\haut{0.16}\def\xg{0.262}
    \pgfmathsetmacro{\hl}{\rE-\rD}% hauteur d'une ligne
    \def\yvc{\haut}\def\yfz{0}

    % V_c et f_z : présents dans les trois parties
    \rpempanneau[expo-]{\cA}{\cB}{\rD}{\rE}{0}{\yvc}{1}
    \rpempanneau[expo-]{\cB}{\cC}{\rD}{\rE}{0}{\yfz}{1}

    % partie 1 : coefficients du modèle a_w ... f_w
    \only<1>{%
        \rpempanneau[expo-]{\cC}{\cD}{\rD}{\rE}{\xg}{\haut}{2}                % a_w
        \rpempanneau{\cD}{\cE}{\rD}{\rE}{\xg+\pas}{\haut}{2}                  % b_w
        \rpempanneau[expo+, expo-]{\cA}{\cB}{\rC}{\rD}{\xg+2*\pas}{\haut}{2}  % c_w
        \rpempanneau[expo+, expo-]{\cB}{\cC}{\rC}{\rD}{\xg}{0}{2}             % d_w
        \rpempanneau[expo+]{\cC}{\cD}{\rC}{\rD}{\xg+\pas}{0}{1}               % e_w
        \rpempanneau[expo-]{\cD}{\cE}{\rC}{\rD}{\xg+2*\pas}{0}{2}}            % f_w

    % partie 2 : pressions spécifiques (ligne 1) et efforts associés (ligne 2)
    \only<2>{%
        \rpempanneau[expo+, expo-, bande=0.039, etiq={$K_f$ (\si{\mega\pascal})}, etiqx=0.014]{\cA}{\cB}{\rB}{\rC}{\xg}{\haut}{0}        % K_f (non exploitable)
        \rpempanneau[expo+, expo-, bande=0.025, etiq={$K_c$ (\si{\mega\pascal})}, etiqx=0.013]{\cB}{0.4895}{\rB}{\rC}{\xg+\pas}{\haut}{1}   % K_c
        \rpempanneau[bande=0.0548, bandey=0.0152, bandeb=0.035, marge=0.006, etiqx=0.0055,
            etiq={$K_m$ (\si{\newton\meter\per\milli\meter\squared})},
            apres={\foreach \yt/\vt in {0.1325/0, 0.0938/$-$2000, 0.0551/$-$4000, 0.0165/$-$6000}
                       \node[anchor=base east, inner sep=0, font=\fontsize{3.5pt}{4pt}\selectfont, text=jvtexte] at (0.053,\yt) {\vt};}
            ]{\cC}{\cD}{\rB}{\rC}{\xg+2*\pas}{\haut}{0}   % K_m (non exploitable)
        \rpempanneau[expo+, bande=0.031, etiq={$F_{\vec{e}_r,\text{max}}$ (\si{\newton})}, etiqx=0.015]{\cA}{\cB}{\rA}{\rB}{\xg}{0}{0}                   % F_er
        \rpempanneau[expo+, bande=0.026, etiq={$F_{\vec{e}_\theta,\text{max}}$ (\si{\newton})}, etiqx=0.012]{\cB}{\cC}{\rA}{\rB}{\xg+\pas}{0}{1}              % F_etheta
        \rpempanneau[bande=0.028, marge=0.012, etiq={$M_{P,\vec{e}_z,\text{max}}$ (\si{\newton\meter})}, etiqx=0.0105]{\cC}{\cD}{\rA}{\rB}{\xg+2*\pas}{0}{0}}   % M_P,ez (non exploitable)

    % partie 3 : énergie spécifique maximale et critère J
    \only<3>{%
        \rpempanneau[expo+, expo-]{\cD}{\cE}{\rB}{\rC}{\xg}{\haut}{1}     % W_c,max
        \rpempanneau[expo+]{\cD}{\cE}{\rA}{\rB}{\xg}{0}{1}}                 % J

    % largeur constante d'une partie à l'autre
    \path (0,\yfz) rectangle (\xg+3*\pas, \haut+\hl);
    % légende des filtres, sous f_z
    \node[anchor=north west, font=\tiny, text=jvtexte, inner sep=0pt]
        at (0, \yfz-0.012)
        {\rpemlegende{1}{1}{1}};
\end{tikzpicture}
